\clearpage\fancyhead[L]{\small PALLAS.}\hypertarget{native-374-lower}{}\label{native-374-lower}
\begin{center}\textbf{PALLAS.}\par\vspace{4pt}\hrule\vspace{4pt}\fontsize{10}{12}\selectfont Göttingische gelehrte Anzeigen. Stück 176. Seite 1753\;..\;1755. 4. November 1813.\par\vspace{4pt}\hrule\vspace{4pt}\end{center}

\par\indent Die Beobachtungen des vergangenen Jahres bestätigten sehr gut die aus den Störungsrechnungen vorausberechnete Bewegung der \emph{Pallas}. Wir haben damals im 67. und 127. Stück dieser Blätter [S. 349 und 356 dieses Bandes] darüber berichtet. Diese Bestätigung erhöhte das Interesse, mit dem die Astronomen die Wiederkehr dieses Planeten im laufenden Jahr erwarteten. An der hiesigen Sternwarte wurde der Planet am 28. Juni zum ersten Mal wieder beobachtet, genau an der Stelle, die Herr \textsc{Nicolai} schon zwei Jahre zuvor in seiner Ephemeride berechnet hatte. Um die Zeit der Opposition. suchte Professor \textsc{Gauss} so viele Beobachtungen wie möglich zu erhalten. Das außergewöhnlich ungünstige Wetter des vergangenen Sommers erlaubte jedoch nur eine geringe Ausbeute. Trotzdem sind diese Beobachtungen vorzüglich ausgefallen, obwohl sie bei dem schwachen Licht des Planeten, der auch diesmal nur von der 10. Größe war, nur mit dem Kreismikrometer angestellt wurden. Hier folgen alle Beobachtungen von Professor \textsc{Gauss}:

\begin{center}\fontsize{10}{12}\selectfont\setlength{\tabcolsep}{5pt}\renewcommand{\arraystretch}{1.15}\begin{tabular}{l|c|c|c}
1813 & Mittlere Zeit & Rektaszension & Nördl. Deklination\\\hline
Juni 28 & $11^{\mathrm u}\;37^{\mathrm m}\;2^{\mathrm s}$ & $327^\circ\;2'\;13''3$ & $14^\circ\;5'\;43''4$\\
August 16 & $10\;36\;23$ & $319\;50\;54.7$ & $.\;.\;.\;.$\\
19 & $9\;31\;0$ & $319\;17\;10.6$ & $10\;0\;8.4$\\
26 & $9\;23\;7$ & $317\;59\;52.0$ & $8\;44\;35.3$\\
\end{tabular}\end{center}

\Needspace{18\baselineskip}\par\indent Herr \textsc{Burckhardt} und Herr \textsc{von Lindenau} teilten Professor \textsc{Gauss} außerdem folgende wertvolle Beobachtungen mit:

\begin{center}Meridianbeobachtungen der Pallas auf der Sternwarte der Militärschule in Paris.\end{center}
\begin{center}\fontsize{10}{12}\selectfont\setlength{\tabcolsep}{5pt}\renewcommand{\arraystretch}{1.15}\begin{tabular}{l|c|c|c}
1813 & Sternzeit & Rektaszension & Nördl. Deklination\\\hline
August 4 & $21^{\mathrm u}\;28^{\mathrm m}\;30^{\mathrm s}5$ & $322^\circ\;7'\;38''0$ & $12^\circ\;11'\;58''5$\\
9 & $21\;24\;42$ & $.\;.\;.\;.$ & $11\;32\;38.4$\\
11 & $.\;.\;.\;.$ & $.\;.\;.\;.$ & $11\;15\;28.1$\\
15 & $21\;20\;7.2$ & $320\;1\;48.0$ & $10\;38\;48.2$\\
17 & $21\;18\;35.2$ & $319\;38\;48.5$ & $10\;19\;29.3$\\
19 & $21\;17\;4.3$ & $319\;16\;4.5$ & $9\;59\;18.3$\\
26 & $21\;11\;56.1$ & $317\;59\;1.8$ & $8\;43\;23.4$\\
\end{tabular}\end{center}

\Needspace{11\baselineskip}\begin{center}\leavevmode\SourceLeaf{375}Beobachtungen der Pallas auf der Seeberger Sternwarte:\end{center}

\begin{center}\fontsize{10}{12}\selectfont\setlength{\tabcolsep}{5pt}\renewcommand{\arraystretch}{1.15}\begin{tabular}{l|c|c}
1813 & Mittlere Zeit & Rektaszension\\\hline
Sept. 1 & $10^{\mathrm u}\;25^{\mathrm m}\;16^{\mathrm s}6$ & $316^\circ\;58'\;15''9$\\
3 & $10\;16\;9.6$ & $316\;39\;25.5$\\
4 & $10\;11\;35.4$ & $316\;29\;48.7$\\
\end{tabular}\end{center}

\Needspace{25\baselineskip}\par\indent Professor \textsc{Gauss} übergab diese Beobachtungen Herrn \textsc{Nicolai}. Dieser geschickte junge Astronom, der nun als Gehilfe von Herrn \textsc{von Lindenau} an der Seeberger Sternwarte angestellt ist, verglich sie zunächst sehr sorgfältig mit den zuletzt verbesserten Elementen und fand die folgende schöne Übereinstimmung:

\begin{center}Abweichung von der Rechnung:\end{center}
\begin{center}\fontsize{10}{12}\selectfont\setlength{\tabcolsep}{7pt}\renewcommand{\arraystretch}{1.15}\begin{tabular}{l|c|c|l}
1813 & Rektaszension & Deklination & Beobachter\\\hline
Juni 28 & $+23''6$ & $+22''4$ & Gauss\\
Aug. 4 & $+36.9$ & $-5.9$ & Burckhardt\\
9 & $.\;.\;.\;.$ & $+5.1$ & Burckhardt\\
11 & $.\;.\;.\;.$ & $+6.2$ & Burckhardt\\
15 & $+26.9$ & $+3.7$ & Burckhardt\\
16 & $+35.6$ & $.\;.\;.\;.$ & Gauss\\
17 & $+28.2$ & $-6.6$ & Burckhardt\\
19 & $+28.5$ & $+4.5$ & Gauss\\
19 & $+24.8$ & $-7.4$ & Burckhardt\\
26 & $+22.4$ & $-4.3$ & Gauss\\
26 & $+18.3$ & $+11.1$ & Burckhardt\\
Sept. 1 & $+22.8$ & $.\;.\;.\;.$ & v. Lindenau\\
3 & $+20.9$ & $.\;.\;.\;.$ & v. Lindenau\\
4 & $+49.9$ & $.\;.\;.\;.$ & v. Lindenau\\
\end{tabular}\end{center}

\Needspace{13\baselineskip}\par\indent Aus den Beobachtungen vom 15., 17. und 19. August leitete Herr \textsc{Nicolai} folgendes Ergebnis für die Opposition ab:

\begin{center}\emph{Neunte beobachtete Opposition der Pallas.}\end{center}
\begin{center}1813 August 18. $8^{\mathrm u}\;41^{\mathrm m}\;5^{\mathrm s}$ mittlere Zeit in Göttingen\end{center}
\begin{center}\fontsize{10}{12}\selectfont\setlength{\tabcolsep}{3pt}\renewcommand{\arraystretch}{1.15}\begin{tabular}{lcr}
Wahre Länge der Pallas & $.\;.\;.\;.\;.\;.\;.\;.$ & $325^\circ\;23'\;56''5$\\
Wahre geozentrische Breite & $.\;.\;.\;.\;.\;.\;.$ & $24\;37\;36.1\ \mathrm{N.}$\\
\end{tabular}\end{center}

\noindent Man vergleiche damit die Vorausberechnung im 127. Stück dieser Anzeigen vom vergangenen Jahr [S. 356 dieses Bandes].

\par\indent Eine erneute Verbesserung der Elemente, um sie auch dieser Opposition noch besser anzupassen, schien Professor \textsc{Gauss} unter diesen Umständen die Mühe nicht wert zu sein.
